\section*{Bab \ref{ch:g8:intensity-ammeter} --- Kuat Arus dan Amperemeter}

\begin{solution}{exo:g8:intensity-ammeter:1}
Laju aliran arusnya --- muatan yang lewat tiap sekon, sebagaimana
sebuah sungai menghitung air tiap sekon. \omterm{def:g6:measurement-in-science:measuring}{Satuannya}: \omterm{def:g8:intensity-ammeter:intensity}{ampere}, \unit{A};
$\qty{250}{mA} = \qty{0.25}{A}$.
\end{solution}

\begin{solution}{exo:g8:intensity-ammeter:2}
\qty{0.02}{A}: lampu penunjuk. \qty{0.3}{A}: \omterm{def:g3:first-electric-circuit:bulb}{bohlam} senter.
\qty{10}{A}: cerek. \qty{100}{A}: mesin yang sedang dinyalakan.
\end{solution}

\begin{solution}{exo:g8:intensity-ammeter:3}
Ia hanya menghitung arak-arakan yang lewat menembus dirinya sendiri,
sehingga ia harus berdiri di tengah jalan --- \omterm{def:g4:series-circuits:series}{terangkai seri}. Ia
hampir sama sekali tidak melawan arusnya; kalau dipasang langsung
melintang pada sebuah \omterm{def:g3:first-electric-circuit:battery}{baterai} ia karena itu menjadi \omterm{def:g7:short-circuits-safety:device}{korsleting}
berpapan ukur.
\end{solution}

\begin{solution}{exo:g8:intensity-ammeter:4}
Matikan lalu putuslah lingkarnya; sisipkan alatnya dengan arus yang
masuk lewat kutub \unit{A} dan keluar lewat kutub bersamanya; mulailah
pada jangkauan yang terbesar lalu turunlah; nyalakan, bacalah, tulislah
\omterm{def:g6:measurement-in-science:measuring}{satuannya}. Jangkauan terbesar lebih dahulu, supaya arus yang tak
terduga kuatnya bertemu setelan yang paling kekar alih-alih merusak
yang halus.
\end{solution}

\begin{solution}{exo:g8:intensity-ammeter:5}
Pada \omterm{def:g4:series-circuits:series}{rangkaian seri} kuat arusnya sama pada setiap titik lingkarnya. Ia
mengubah menjadi bilangan pengamatan lama bahwa dua \omterm{def:g3:first-electric-circuit:bulb}{bohlam} seri yang
serupa menyala persis sama terangnya, di mana pun mereka duduk.
\end{solution}

\begin{solution}{exo:g8:intensity-ammeter:6}
\qty{0.45}{A}, dan \qty{0.45}{A}: satu lingkar, satu \omterm{def:g8:intensity-ammeter:intensity}{kuat arus}.
\end{solution}

\begin{solution}{exo:g8:intensity-ammeter:7}
Tidak. Hukumnya menjanjikan satu nilai \emph{mengelilingi lingkarnya
pada saat tertentu}, bukan nilai yang sama sesudah lingkarnya dibangun
ulang. Lampu keduanya melemahkan arak-arakannya --- sama rata di
mana-mana: semua gardu tolnya kini sepakat pada \qty{0.17}{A}.
\end{solution}

\begin{solution}{exo:g8:intensity-ammeter:8}
Memasang alatnya terbalik --- arus masuk lewat kutub bersamanya. Pada
alat digital, tidak berbahaya: tanda minusnya adalah seluruh hukumannya
(alat berjarum akan menghantam penahannya).
\end{solution}

\begin{solution}{exo:g8:intensity-ammeter:9}
$0.60 - 0.35 = \qty{0.25}{A}$ --- percabangannya harus seimbang: apa
yang mengalir masuk mengalir keluar. Gardu tol pada percabangan yang
kedua membaca \qty{0.60}{A} yang bersatu kembali.
\end{solution}

\begin{solution}{exo:g8:intensity-ammeter:10}
Ketika terpasang paralel, alatnya yang hampir tanpa upaya menjadi
sebuah jalan pintas: lampunya, karena ditinggalkan, meredup sampai
tiada, sementara alatnya membawa arus yang menyerbu yang dibatasi oleh
hampir tidak ada --- sebuah \omterm{def:g7:short-circuits-safety:device}{korsleting} menembus alatnya, yang
memutuskan \omterm{def:g7:short-circuits-safety:fuse}{sekring} di dalamnya kalau nasibnya baik. Muridnya membangun
penjahat bab keselamatannya.
\end{solution}

\begin{solution}{exo:g8:intensity-ammeter:11}
Jalan utamanya membawa jumlahnya: \qty{0.60}{A}. \omterm{def:g5:series-parallel-circuits:parallel}{Cabang} ketiga yang
serupa membawanya ke \qty{0.90}{A} --- dan \omterm{def:g3:first-electric-circuit:battery}{baterainya}, yang melayani
tiga arak-arakan penuh, terkuras paling cepat di antara semuanya.
\end{solution}

\begin{solution}{exo:g8:intensity-ammeter:12}
Arus yang harus diketahui: $I$ jalan utamanya, dan arus tiap \omterm{def:g5:series-parallel-circuits:parallel}{cabangnya}
($I_1$, $I_2$) --- tiga seluruhnya (ruas lingkar di luar percabangannya
mengulangi ketiganya). Dua pengukuran mencukupi: ukurlah dua yang mana
pun, dan keseimbangan percabangan $I = I_1 + I_2$ menyerahkan yang
ketiga. Urutan yang aman: rangkaiannya dimatikan; sisipkan alatnya ke
jalan yang dipilih; jangkauan terbesar; nyalakan; baca; matikan
sebelum memindahkan alatnya --- dan alatnya tidak pernah, pada langkah
mana pun, menjembatani dua titik pada jalan yang berbeda.
\end{solution}

\begin{solution}{pb:g8:intensity-ammeter:1}
\textbf{1.} Matikan; putuskan lingkarnya di antara lampu dan
\omterm{def:g3:first-electric-circuit:battery}{baterainya}; sisipkan \omterm{def:g8:intensity-ammeter:ammeter}{amperemeternya}, arus masuk lewat \unit{A}, keluar
lewat kutub bersamanya; jangkauan terbesar, lalu turun; nyalakan;
baca; tulis \omterm{def:g6:measurement-in-science:measuring}{satuannya}.
\textbf{2.} \qty{0.28}{A}.
\textbf{3.} \qty{280}{mA} lagi --- satu lingkar, satu \omterm{def:g8:intensity-ammeter:intensity}{kuat arus}:
gardu tolnya menghitung arak-arakan yang sama di kedua sisinya.
\textbf{4.} Di mana pun pada lingkarnya --- keduanya \omterm{def:g4:series-circuits:series}{terangkai seri},
pada kedudukan mana pun: keduanya menunjukkan \qty{280}{mA}.
\textbf{5.} Hukumnya menolak ramalannya: lampu seri membawa arus yang
sama persis. Perbedaannya terletak pada lampunya sendiri --- bangunan
yang tidak sama mengubah arak-arakan yang sama menjadi \omterm{def:g1:light-and-shadows:source}{cahaya} yang
tidak sama (bab berikutnya menamai apa yang berbeda).
\textbf{6.} Setiap titik lingkarnya: \qty{150}{mA} --- sebelum A,
sesudah B, dan di samping \omterm{def:g3:first-electric-circuit:battery}{baterainya}.
\textbf{7.} B mati (ditinggalkan oleh jalan pintasnya) dan lingkarnya,
yang telah kehilangan satu penentang, membawa arak-arakan yang lebih
kuat: alatnya memanjat --- dan lampu A menjadi lebih terang sesuai
dengan itu.
\textbf{8.} ``Melintang pada B jalan pintasnya masih meninggalkan
lampu A yang menjaga lingkarnya; melintang pada \omterm{def:g3:first-electric-circuit:battery}{baterainya} ia sama
sekali tidak meninggalkan apa pun --- dan arak-arakan yang tak terjaga
adalah bab kebakaran, bukan sebuah peragaan.''
\textbf{9.} $0.52 - 0.30 = \qty{0.22}{A}$.
\textbf{10.} $0.40 + 0.30 = \qty{0.70}{A}$ --- dan \omterm{def:g7:short-circuits-safety:fuse}{sekring} \omterm{def:g5:series-parallel-circuits:parallel}{cabang}
(atau pemutus arus) yang menjaga jalur \omterm{ex:g6:circuits-and-safety:motor}{motornya} mendengarkan dengan
saksama: \omterm{ex:g6:circuits-and-safety:motor}{motor} yang macet adalah cara pengawal memperoleh nafkahnya.
\textbf{11.} \omterm{def:g5:series-parallel-circuits:parallel}{Cabang} lampunya: \qty{0}{A}; \omterm{def:g5:series-parallel-circuits:parallel}{cabang} \omterm{ex:g6:circuits-and-safety:motor}{motornya}:
\qty{0.40}{A} miliknya sendiri, tidak berubah (kemandirian paralel);
jalan utamanya: \qty{0.40}{A}.
\textbf{12.} Kemurahan hatinya dibayar di \omterm{def:g3:first-electric-circuit:battery}{baterainya}: dua \omterm{def:g5:series-parallel-circuits:parallel}{cabang} yang
sehat meminum dua arak-arakan penuh, dan arus jalan utama yang
terjumlah mengosongkan simpanannya paling cepat.
\textbf{13.} Misalnya: ``Sebuah \omterm{def:g8:intensity-ammeter:ammeter}{amperemeter} berdiri di tengah jalan,
tidak pernah melintang padanya. Satu lingkar menjamin satu \omterm{def:g8:intensity-ammeter:intensity}{kuat arus}
--- di mana-mana, selalu. Dan sebuah percabangan menyeimbangkan
pembukuannya: arak-arakan masuk sama dengan arak-arakan keluar, sampai
ke miliamperenya.''
\end{solution}
